\documentclass[11pt]{article}
\usepackage[margin=0.85in]{geometry}
\usepackage{amsmath,amstext,amssymb,amsthm}
\usepackage{multirow}
\usepackage{tikz,pgfplots}
\pgfplotsset{compat=1.18}
\usepackage{graphicx}
\graphicspath{{./img/}}
\usepackage{caption,subcaption}
\usepackage{xcolor}
\usepackage{tcolorbox}
\usepackage{hyperref}

\hypersetup{
    colorlinks=true,
    linkcolor=blue,
    urlcolor=blue
}
\setlength{\parindent}{0in}
\setlength{\parskip}{0.5em}

\providecommand{\la}{\langle}
\providecommand{\ra}{\rangle}
\providecommand{\myvec}[1]{\mathbf{#1}}

% Render \content{main}{presentation}{nopresentation}{comments} with answers/notes visible
\newcommand{\content}[4]{%
    #1\par
    \if\relax\detokenize{#3}\relax\else
        #3\par
    \fi
    \if\relax\detokenize{#4}\relax\else
        \begin{tcolorbox}[colback=blue!5!white, colframe=blue!40!black, boxrule=0.5pt, arc=2pt, left=6pt, right=6pt, top=4pt, bottom=4pt]
            \small\textbf{Notes / Answer:} #4
        \end{tcolorbox}
    \fi
    \medskip
}

% Common math & theorem environments
\providecommand{\pd}[2]{\frac{\partial #1}{\partial #2}}
\providecommand{\R}{\mathbb{R}}
\providecommand{\N}{\mathbb{N}}
\providecommand{\Z}{\mathbb{Z}}
\providecommand{\ve}[1]{\mathbf{#1}}

\theoremstyle{definition}
\newtheorem{theorem}{Theorem}[section]
\newtheorem{definition}[theorem]{Definition}
\newtheorem{example}[theorem]{Example}
\newtheorem{remark}[theorem]{Remark}
\newtheorem{note}[theorem]{Note}
\newtheorem{proposition}[theorem]{Proposition}
\newtheorem{corollary}[theorem]{Corollary}
\newtheorem{lemma}[theorem]{Lemma}

\begin{document}

\title{Calculus 3 Lecture Notes}
\author{Braden Carlson}
\date{}
\maketitle
\tableofcontents
\newpage

\section{Three Dimensional Coordinate Systems and Vectors}
\input{lecture-1.tex}
\input{lecture-2.tex}
\input{lecture-3.tex}
\input{lecture-4.tex}
\input{lecture-5.tex}
\input{lecture-6.tex}
\input{lecture-7.tex}

\section{Vector Functions}
\input{lecture-8.tex}
\input{lecture-9.tex}
\input{lecture-10.tex}

\section{Partial Derivatives}
\input{lecture-11.tex}
\input{lecture-12.tex}
\input{lecture-13.tex}
\input{lecture-14.tex}
\input{lecture-15.tex}
\input{lecture-16.tex}
\input{lecture-17.tex}
\input{lecture-18.tex}

\section{Multiple Integrals}
\input{lecture-19.tex}
\input{lecture-20.tex}
\input{lecture-21.tex}
\input{lecture-22.tex}
\input{lecture-23.tex}
\input{lecture-24.tex}
\input{lecture-25.tex}
\input{lecture-26.tex}
\input{lecture-27.tex}

\section{Vector Calculus}
\input{lecture-28.tex}
\input{lecture-29.tex}
\input{lecture-30.tex}
\input{lecture-31.tex}
\input{lecture-32.tex}
\input{lecture-33.tex}
\input{lecture-34.tex}
\input{lecture-35.tex}
\input{lecture-36.tex}

\end{document}
